\section{练习}

\begin{question}\begin{drill}[同号双电荷中垂线上的最大场]\label{e1:drill:1}
两相同正点电荷 $q$ 相距 $2a$。求其中垂线上电场大小的最大值及位置。
\end{drill}
\typeout{ECFIG:fig:e1:drill:1:s}\fig{fig:e1:drill:1:s}{题设：同号双电荷中垂线上的最大场}{\begin{tikzpicture}[scale=.88]
\draw[world,fill=white](-1,0)circle(.12);\node[above]at(-1,0.15){$q$};\draw[world,fill=white](1,0)circle(.12);\node[above]at(1,0.15){$q$};\draw[axis](0,-.3)--(0,3)node[above]{$y$};\fill(0,2)circle(1.5pt)node[right]{$P$};\node at(0,-.7){间距 $2a$};
\end{tikzpicture}}
\end{question}
\answerlink{e1:drill:1}
\begin{question}\begin{drill}[绝缘球环中的轴向小振动]\label{e1:drill:2}
均匀绝缘圆环半径 $R$、总电荷 $Q>0$。质量 $m$ 的点电荷 $q$ 被约束只沿垂直环面的轴线运动。要使它在环心附近小振动，$q$ 应取什么符号？求周期与小振动条件。
\end{drill}
\typeout{ECFIG:fig:e1:drill:2:s}\fig{fig:e1:drill:2:s}{题设：绝缘球环中的轴向小振动}{\begin{tikzpicture}[scale=.88]
\draw[world](0,0)ellipse(.5 and 1.4);\draw[axis](-1,0)--(4.5,0)node[right]{$x$};\fill(0,0)circle(1.3pt)node[below left]{$O$};\fill(3,0)circle(1.6pt)node[above]{$P$};\draw[dim](0,-1.8)--(3,-1.8)node[midway,below]{$x$};\draw[dim](0,0)--(0,1.4)node[midway,left]{$R$};
\end{tikzpicture}}
\end{question}
\answerlink{e1:drill:2}
\begin{question}\begin{drill}[圆环外两个点电荷同时平衡]\label{e1:drill:3}
圆环半径 $R$，固定均匀带电 $Q>0$。轴线上 $x=r>R$ 放非零电荷 $q_1$，问能否在 $x>0$ 的另一位置放非零电荷 $q_2$，使两个点电荷都平衡？求位置及电荷关系。
\end{drill}
\typeout{ECFIG:fig:e1:drill:3:s}\fig{fig:e1:drill:3:s}{圆环正半轴上的内、外两类候选位置；固定的 $q_1$ 坐标均为 $r$}{\begin{tikzpicture}

\begin{scope}
\draw(0,0)ellipse(.35 and 1.2);\draw[axis](-.8,0)--(5,0)node[right]{$x$};
\fill(2,0)circle(2pt)node[above]{$q_2$};\fill(4,0)circle(2pt)node[above]{$q_1$};
\node[below]at(2,0){$x$};\node[below]at(4,0){$r$};\node at(2,-1.7){内侧候选 $0<x<r$};\node[above]at(0,1.3){$Q>0,\ R$};
\end{scope}
\begin{scope}[xshift=7.5cm]
\draw(0,0)ellipse(.35 and 1.2);\draw[axis](-.8,0)--(5,0)node[right]{$x$};
\fill(2,0)circle(2pt)node[above]{$q_1$};\fill(4,0)circle(2pt)node[above]{$q_2$};
\node[below]at(2,0){$r$};\node[below]at(4,0){$x$};\node at(2,-1.7){外侧候选 $x>r$};\node[above]at(0,1.3){$Q>0,\ R$};
\end{scope}

\end{tikzpicture}}
\end{question}
\answerlink{e1:drill:3}
\begin{question}\begin{drill}[导体表面的电压力]\label{e1:drill:4}
静电平衡导体表面外侧法向场强为 $E$。求表面电荷受到的向外电压力，并说明为什么有效受力场只有外侧总场的一半。
\end{drill}
\typeout{ECFIG:fig:e1:drill:4:s}\fig{fig:e1:drill:4:s}{题设：导体表面的电压力}{\begin{tikzpicture}[scale=.88]
\fill[gray!15](-2,-1)rectangle(2,0);\draw[world](-2,0)--(2,0);\draw[flow](0,.2)--(0,1.7)node[right]{$E$};\node at(0,-.5){$E_{\text{内}}=0$};
\end{tikzpicture}}
\end{question}
\answerlink{e1:drill:4}
\begin{question}\begin{drill}[均匀有限直线的中垂面电场]\label{e1:drill:5}
均匀带电直线段长 $2l$、线密度 $\lambda>0$，位于 $y$ 轴 $-l\le y\le l$。求 $P=(x,0),x>0$ 处电场，并讨论 $l\to\infty$ 及 $x\gg l$。
\end{drill}
\typeout{ECFIG:fig:e1:drill:5:s}\fig{fig:e1:drill:5:s}{题设：均匀有限直线的中垂面电场}{\begin{tikzpicture}[scale=.88]
\draw[axis](-.2,0)--(4.2,0)node[right]{$x$};\draw[axis](0,-.2)--(0,2.5)node[above]{$y$};\draw[world,very thick](0,-2)--(0,2);\node[left]at(0,2){$l$};\node[left]at(0,-2){$-l$};\fill(3,0)circle(1.7pt)node[above]{$P$};
\end{tikzpicture}}
\end{question}
\answerlink{e1:drill:5}
\begin{question}\begin{drill}[带电圆弧在圆心的电场]\label{e1:drill:6}
半径 $R$ 的圆弧关于 $x$ 轴对称，覆盖极角 $-\alpha\le\theta\le\alpha$，$0<\alpha\le\pi$，线密度 $\lambda>0$。求圆心场，并给出半圆弧情形。
\end{drill}
\typeout{ECFIG:fig:e1:drill:6:s}\fig{fig:e1:drill:6:s}{题设：带电圆弧在圆心的电场}{\begin{tikzpicture}[scale=.88]
\draw[axis](-1,0)--(3,0)node[right]{$x$};\draw[world](1.026,-2.819)arc(-70:70:3);\fill(0,0)circle(1.4pt)node[below]{$O$};\draw[guide](0,0)--(1.026,2.819)(0,0)--(1.026,-2.819);\node at(1.1,.7){$\alpha$};\node at(1.1,-.7){$\alpha$};\node at(2.9,1.3){$R$};
\end{tikzpicture}}
\end{question}
\answerlink{e1:drill:6}
\begin{question}\begin{drill}[无限均匀带电圆柱的内外电场]\label{e1:drill:7}
无限长绝缘圆柱半径 $R$、体电荷密度 $\rho>0$。求距轴 $r$ 处电场，并检查表面连续性。
\end{drill}
\typeout{ECFIG:fig:e1:drill:7:s}\fig{fig:e1:drill:7:s}{题设：无限均匀带电圆柱的内外电场}{\begin{tikzpicture}[scale=.88]
\draw[world](0,0)ellipse(.8 and 1.4);\draw[world](0,1.4)--(4,1.4)(0,-1.4)--(4,-1.4);\draw[world](4,0)ellipse(.8 and 1.4);\draw[guide](-1,0)--(5,0);\node at(2,1.9){$\rho$};\node at(-1,1){$R$};
\end{tikzpicture}}
\end{question}
\answerlink{e1:drill:7}
\begin{question}\begin{drill}[均匀实心球的内外场与最大值]\label{e1:drill:8}
绝缘球半径 $R$，均匀体电荷密度 $\rho>0$。求内外电场，判断场强最大位置；若半径变为 $2R$，分别讨论保持 $\rho$ 与保持总电荷的变化。
\end{drill}
\typeout{ECFIG:fig:e1:drill:8:s}\fig{fig:e1:drill:8:s}{题设：均匀实心球的内外场与最大值}{\begin{tikzpicture}[scale=.88]
\draw[world](0,0)circle(1.5);\fill(0,0)circle(1.2pt)node[below]{$O$};\draw[dim](0,0)--(-1.06,1.06)node[midway,left]{$R$};\node at(2.4,.5){$\rho$ 均匀};
\end{tikzpicture}}
\end{question}
\answerlink{e1:drill:8}
